% Round 18. ci_global branch. Independent review required.
\section{Ample-cut curve comparison and surface ambient Green families}
\label{r18sc:section}

PG18's numerical $U_2$ theorem does not require the whole ambient
$H^2$ to be algebraic.  Its separate representative gate can also
be closed by a curve-comparison property of the ambient variety.
The new closure operation in this section is propagation of that
property through smooth ample cuts of final dimension at least
three.  The projective-bundle and smooth-blowup curve comparisons
themselves are already frozen in CD17 and are not reopened.

\subsection{The native curve property and its ample-cut propagation}

Let $k$ be a finite field, fix $\ell\ne\operatorname{char}k$,
and let $T/k$ be smooth projective geometrically integral.
Say that $T$ has the curve comparison property $\mathrm{CC}_\ell$
if every actual native rational curve cycle which is numerically
zero has zero class in
\[
 H^{2\dim T-2}(T_{\bar k},\mathbf Q_\ell(\dim T-1)).
 \tag{SC18.1}\label{r18sc:CC}
\]
This is a statement about actual cycles and their geometric
classes.  It does not start by postulating a class map on
an arbitrary numerical quotient.

\begin{lemma}[Smooth ample cuts preserve curve comparison]
\label{r18sc:ample-cuts}
Let
\[
 Y=T_0\supset T_1\supset\cdots\supset T_c=E
 \tag{SC18.2}\label{r18sc:cuts}
\]
be native smooth projective geometrically integral varieties,
each $T_j$ a smooth ample Cartier divisor in $T_{j-1}$.
Suppose $\dim E\ge3$.  If $Y$ has $\mathrm{CC}_\ell$,
then $E$ has $\mathrm{CC}_\ell$.
If the ambient property holds afresh after every finite
extension of $k$, the conclusion does too.
\end{lemma}

\begin{proof}
Write $m=\dim Y$, $e=\dim E$, so $c=m-e$, and let
$i:E\hookrightarrow Y$ be the actual closed immersion.
At each cut the target dimension is at least three,
so geometric etale weak Lefschetz gives an isomorphism
in degree two.  Their composite is therefore
\[
 i^*:H^2(Y_{\bar k},\mathbf Q_\ell(1))
               \xrightarrow{\sim}H^2(E_{\bar k},\mathbf Q_\ell(1)).
 \tag{SC18.3}\label{r18sc:H2iso}
\]
Poincare duality and the actual cohomological Gysin
adjunction identify the dual map as
\[
 i_*:H^{2e-2}(E_{\bar k},\mathbf Q_\ell(e-1))
               \xrightarrow{\sim}
          H^{2m-2}(Y_{\bar k},\mathbf Q_\ell(m-1)).
 \tag{SC18.4}\label{r18sc:Gysin-iso}
\]
In particular this particular Gysin map is injective;
the injection has been proved by duality, not inferred
by cancelling an ambient pushforward equality.

Let $z$ be an actual native numerically zero rational
curve on $E$.  For every native divisor $L$ on $Y$,
the proper degree projection formula gives
\[
                  \deg_Y(i_*z\cdot L)=\deg_E(z\cdot i^*L)=0.
\]
Thus $i_*z$ is numerically zero as a curve on $Y$.
The ambient hypothesis gives
$\operatorname{cl}_\ell(i_*z)=0$.
Cycle-class functoriality and the established injection
\eqref{r18sc:Gysin-iso} force
$\operatorname{cl}_\ell(z)=0$.
All degrees were native degrees on the proper strata.
The same proof works after any finite residue extension
to which the ambient hypothesis is reapplied.
\end{proof}

\begin{lemma}[The edge blowups introduce only zero-cycle coefficients]
\label{r18sc:edge-blowups}
If $E$ has $\mathrm{CC}_\ell$, then its specified
codimension-two smooth blowups in geometrically
integral centers have $\mathrm{CC}_\ell$ as well.
\end{lemma}

\begin{proof}
This is the relevant case of frozen CD17's
smooth-modification theorem.  An explicit coefficient
check is useful here.  If $d=\dim E$,
$q:\widetilde E=\operatorname{Bl}_ZE\to E$,
$j:J\hookrightarrow\widetilde E$ and $p:J\to Z$,
the actual Chow decomposition in curve degree is
\[
 z=q^*z_0+j_*p^*\beta,
 \qquad z_0\in\operatorname{CH}^{d-1}(E)_{\mathbf Q},
 \quad\beta\in\operatorname{CH}^{d-2}(Z)_{\mathbf Q}.
 \tag{SC18.5}\label{r18sc:edge-decomposition}
\]
The second coefficient is a zero-cycle.
Here $q_*z=z_0$, while $\beta=-p_*j^*z$:
$j^*q^*z_0=p^*i^*z_0=0$ by dimension, and
$j^*j_*=-\xi$, $p_*\xi=1$ give the latter formula.
These actual extraction operations preserve numerical
equivalence by the projection formula.  A numerically
zero $z$ therefore has numerically zero $z_0$ and
a degree-zero $\beta$.  The first has zero geometric
class by hypothesis, and the second has zero top
class on the smooth geometrically integral center.
Functoriality applied to the actual decomposition
gives $\operatorname{cl}_\ell(z)=0$.
Iteration proves the assertion.  The proof holds
over every further finite native residue field.
\end{proof}

\begin{theorem}[The ambient curve gate in PG18]
\label{r18sc:PG-gate}
Assume all of PG18's native product-smoothing
hypotheses on a smooth projective ambient
$\mathcal Y/R$ of relative dimension $n+1\ge5$,
with a relatively very ample $\mathcal L$ and
the specified transverse factors and smoothing
section.  If its special fiber $Y$ has
$\mathrm{CC}_\ell$, then every specified
geometrically homologically trivial rational
codimension-two cycle on the generic fiber
has strong G.  In dimension $n=4$ this is
full G in all codimensions.

If the ambient curve property holds after
every finite residue extension, the theorem
applies afresh to all new cycles after every
unramified valued extension.  With RI18 this
gives the corresponding full fourfold assertion
after every finite valued extension.
\end{theorem}

\begin{proof}
Each starting double $H_i\cap H_j$ is obtained
from $Y$ by two smooth ample cuts, with final
dimension $n-1\ge3$.  Lemma~\ref{r18sc:ample-cuts}
gives its curve comparison.  PG18's actual
ordered model makes the final double its
specified codimension-two blowups in smooth
geometrically integral centers of dimension
$n-3$.  Lemma~\ref{r18sc:edge-blowups} gives
the needed comparison on each final double.

Thus the exact curve hypothesis of PG18's
Green corollary is satisfied.  Its numerical
theorem gives $U_2=0$; CL17 gives
$K_{n-2}=B_{n-2}$ from the comparison; DQ17
and ER17 annihilate the specified-cycle
obstruction.  VT17 and frozen reconstruction
give the ordinary actual form on the complement
of the original precise support and the identity
against every original compact ordinary actual
test.  The frozen endpoint and one-cycle
theorems finish the other fourfold degrees.

An unramified extension preserves the native
model equations and geometric conditions.
Reapply the ambient property and the two
lemmas over the new residue field and to
its entire new cycle groups.  RI18's
separate theorem then gives the full
fourfold assertion under arbitrary finite
valued extension.  No claim that those
new cycles are original-field pullbacks
is made.
\end{proof}

\subsection{Good surfaces and their projective bundles}

\begin{corollary}[Projective bundles over arbitrary good surfaces]
\label{r18sc:surface-bundle}
Let $\mathcal S/R$ be smooth projective with
geometrically integral surface fibers, and
let $\mathcal V$ be a rank-$n$ vector bundle
on it, where $n\ge4$.  Set
$\mathcal Y=\mathbf P_{\mathcal S}(\mathcal V)$,
of relative dimension $n+1$.
Every product smoothing on this ambient
with all PG18 hypotheses has codimension-two
strong G.  For $n=4$ it has full G in
every codimension and, using RI18, the
same full conclusion for all cycles over
every finite extension of the valued field.
No algebraicity assumption on the whole
geometric $H^2(\mathcal S_k)$ is required.
\end{corollary}

\begin{proof}
CD17's projective-bundle curve comparison
is uniform after every finite residue
extension.  It applies to the actual
special fiber of this projective bundle.
Theorem~\ref{r18sc:PG-gate} therefore
applies with precisely its stated hypotheses.
\end{proof}

For the product $Y=S\times\mathbf P^{n-1}$
the native curve comparison can be seen
without any bundle coefficient shorthand.
Writing $h$ for the projective-factor class,
every actual rational curve class has the
unique decomposition
\[
 z=p_S^*D\,h^{n-1}+p_S^*z_0\,h^{n-2},
 \quad D\in\operatorname{CH}^1(S)_{\mathbf Q},
 \quad z_0\in\operatorname{CH}^2(S)_{\mathbf Q}.
 \tag{SC18.6}\label{r18sc:product-curves}
\]
If $z$ is numerically zero, testing against
every $p_S^*L$ makes $D$ numerically zero
on the surface.  Testing against $h$ gives
$\deg z_0=0$.  The divisor comparison and
the top-degree class of a zero-cycle make
both geometric coefficient classes zero.
Their pullback products therefore make
$\operatorname{cl}_\ell(z)=0$.
This proof reapplies over every finite
extension of the native finite field and
uses no divisor Tate conjecture for $S$.

\begin{lemma}[Local examples exist after an explicitly allowed unramified extension]
\label{r18sc:local-existence}
For any supplied $\mathcal S/R$ as above,
choose a relatively very ample line bundle
$\mathcal A$ on it and put
\[
 \mathcal Y=\mathcal S\times_R\mathbf P_R^{n-1},
       \qquad\mathcal L=\mathcal A\boxtimes\mathcal O(1).
 \tag{SC18.7}\label{r18sc:product-polarization}
\]
For any $s\ge3$ and positive factor degrees,
there is a finite unramified $R'/R$ over
which actual factor and smoothing sections
satisfying every PG18 hypothesis exist.
The generic field retains the extension
of the original absolute value.
\end{lemma}

\begin{proof}
Use an actual relative projective embedding
given by $\mathcal L$ and the spaces of
ambient homogeneous polynomials of the
specified degrees.  Over $\bar k$, choose
the factor sections successively by Bertini
on the finite collection of prior smooth
strata, avoiding the prior zero-dimensional
strata.  Choose the smoothing section
transverse to all positive-dimensional
factor strata and avoiding their zero-
dimensional intersections.  These are
finitely many nonempty open conditions.
Every resulting positive-dimensional
intersection, including one with the
smoothing section, is geometrically
connected by successive weak Lefschetz
in degree zero, and smoothness makes
it integral.

The nonempty parameter open over $\bar k$
has a point over some finite $k'/k$.
Take the corresponding finite unramified
extension $R'/R$ and lift the ambient
polynomial coefficients.  Lifting those
coefficients uses their free parameter
module; it does not assume surjectivity
of a possibly varying section-cohomology
base-change map.  The special geometric
conditions are now fixed.

Fix lifts of the factors and one smoothing
section $Q_0$.  The other lifts of that
section are $Q_0+\pi T$ in ambient
polynomial coefficients.  The generic
divisor equations are consequently
$\prod_i f_i+\pi Q_0+\pi^2T$.
This affine coefficient coset is Zariski
dense over $\operatorname{Frac}(R')$:
each coefficient ranges over an infinite
arithmetic coset in that infinite field.
It therefore meets the nonempty generic
smooth-section open, by Bertini for the
supplied very ample system.  Such a
smooth ample divisor is geometrically
connected, hence integral.  This fixes
all the hypotheses, including generic
smoothness, without changing the special
data.
\end{proof}

The constructed variety may be defined only
over $\operatorname{Frac}(R')$.  That is
the field over which this existence result
asserts its full cycle domain; it does not
silently descend the variety.  If an
already specified original-field variety
acquires this supplied model over a finite
auxiliary extension, the frozen full-degree
finite actual detector descends G for its
original cycles with the original absolute
value.  This is a separate, already proved
descent use.

\subsection{A concrete example with a nonalgebraic ambient degree two}

\begin{proposition}[The ambient algebraic-$H^2$ condition is unnecessary]
\label{r18sc:elliptic-example}
Let $K=\mathbf Q_5$ and take the good
elliptic curve $E_0:y^2=x^3-x$.
Let $\mathcal S=\mathcal E_0\times_R\mathcal E_0$,
$\mathcal A=\mathcal O_{\mathcal E_0}(3[0])
\boxtimes\mathcal O_{\mathcal E_0}(3[0])$,
and $\mathcal Y=\mathcal S\times\mathbf P_R^3$
with the polarization \eqref{r18sc:product-polarization}.
After a finite unramified extension there are
five-factor, degree-one product smoothings
whose generic fourfold has full G and
\[
             b_2(X)=7,\qquad\operatorname{rank}N^4=179.
 \tag{SC18.8}\label{r18sc:example-invariants}
\]
The ambient geometric $H^2(1)$ is not spanned
by divisor classes, even after further finite
residue extension.  Thus this example uses
the curve-comparison gate rather than the
stronger algebraic-$H^2$ sufficient condition.
\end{proposition}

\begin{proof}
The discriminant of the displayed integral
elliptic equation is $64$, a $5$-adic unit.
Its reduction has eight points over
$\mathbf F_5$: the affine point counts
for $x=0,1,2,3,4$ are respectively
$1,1,2,2,1$, followed by the point
at infinity.  Its Frobenius polynomial
on $H^1$ is therefore $T^2+2T+5$,
with roots $\alpha=-1+2i$ and
$\beta=-1-2i$.
The Kunneth decomposition of the product
surface has two eigenvalues
\[
                    \alpha^2/5=(-3-4i)/5,
                    \qquad\beta^2/5=(-3+4i)/5
\]
on its Tate-twisted degree-two cohomology.
They are not roots of unity: the roots
of unity in $\mathbf Q(i)$ are
$1,-1,i,-i$.  Every geometric divisor
over a finite field is defined over a
finite extension, so some power of
Frobenius fixes its Tate-twisted class.
Divisor classes cannot span these two
eigenspaces.  Raising the eigenvalues
to any positive residue-extension power
does not make them roots of unity.
Adding the projective factor only adds
an algebraic hyperplane line.

Nevertheless the product surface ambient
has the uniform curve comparison
\eqref{r18sc:product-curves}.
The local existence lemma supplies all
five-factor PG18 data after one finite
unramified extension, and
Corollary~\ref{r18sc:surface-bundle}
gives full G for its entire fourfold
cycle group, and then RI18 gives all
finite extensions as well.

On the surface $\deg(A^2)=2\cdot3\cdot3=18$.
On $Y=S\times\mathbf P^3$, write
$h=A+b$.  Its top degree is
\[
                        \deg_Y(h^5)=\binom52\deg_S(A^2)=180.
\]
The fivefold factor stratum therefore
has 180 geometric points, while every
lower positive-dimensional stratum is
connected.  The HM18/NG18 extreme
calculation gives rank $N^4=180-1$.
All points and their native residue
degrees are retained.
Finally weak Lefschetz identifies
$H^2(X)$ with $H^2(Y)$; an elliptic
product surface has $b_2=1+4+1=6$,
and $\mathbf P^3$ adds one, giving
$b_2(X)=7$.  In particular this is
not geometrically a projective-space
complete-intersection fourfold, whose
second Betti number is one.
The positive fourth monodromy also
excludes potential good reduction
and crossing depth less than four
for strict semistable dominant
witnesses, by NM17.
\end{proof}

\subsection{A global theorem when a good surface model is supplied}

\begin{theorem}[Surface ambient full-height families under a supplied global model]
\label{r18sc:global-input}
Let $F$ be a number field and let
$\mathcal S/\mathcal O_F$ be a supplied
smooth projective scheme with geometrically
integral surface fibers.  Fix a relatively
very ample $\mathcal A$ and a relative
projective embedding of
$\mathcal Y=\mathcal S\times\mathbf P^3$
given by $\mathcal L=\mathcal A\boxtimes\mathcal O(1)$.
Then, over one fixed finite extension
$E/F$, there is a Zariski-dense family
in the linear system supplied by degree-five
polynomials of that embedding whose smooth
generic divisors have full local G at every
finite place and complete complementary
homological Chow-group BB heights over
every finite extension of $E$.
At a specified good rational prime,
their product degeneration has
\[
                         \operatorname{rank}N^4=10\deg(A^2)-1>0.
 \tag{SC18.9}\label{r18sc:global-monodromy}
\]
If the supplied degree-five polynomials
span $H^0(Y_E,\mathcal L^5)$, the density
is in that full linear system.
\end{theorem}

\begin{proof}
All fibers of the ambient are smooth
and geometrically integral.  The fixed-
member pencil proof of CG18 uses only
the first and second jets and separated
two-point first jets of a degree at
least three in a projective embedding,
so it applies to this degree-five
section system on the actual ambient.
The special smooth-section opens are
geometrically nonempty by the same
first-jet incidence argument.

Their smooth and one-node members are
geometrically integral for a direct
reason which does not require $S$
to be a projective complete intersection.
For every geometric fiber,
\[
 H^1\bigl(S\times\mathbf P^3,
                  A^{-5}\boxtimes\mathcal O(-5)\bigr)=0
\]
by Kunneth: both possible terms have
either $H^0(\mathbf P^3,\mathcal O(-5))$
or $H^1(\mathbf P^3,\mathcal O(-5))$,
which vanish.  The divisor sequence
therefore gives $H^0(X,\mathcal O_X)=\bar k$.
A smooth or isolated-node Cartier
fourfold is $S_2$ and $R_1$, so it
is connected normal and integral.

Choose a smooth generic reference
section.  Let $S_0$ be the finite set
of rational primes under its bad
places, together with $2$, and select
every place $v$ of $F$ over every
$p\in S_0$, including the places
where the reference happens to be good.
Write $\kappa(v)=\mathbf F_{p^{f_v}}$.
The geometrically nonempty smooth-section
open has a point over some
$\mathbf F_{p^{f_vt_v}}$.
Take $N$ divisible by all $f_vt_v$,
and choose a monic polynomial over
$\mathbf Z$ whose reductions at all
$p\in S_0$ are irreducible of degree
$N$.  It defines a number field $L$
unramified and inert of degree $N$
at these primes.  In $E=FL$, every
residue field above $v$ contains
$\mathbf F_{p^{f_vt_v}}$ and hence
the chosen smooth section.  These
congruences are imposed at every
prime of $E$ over $S_0$.
This saturation is necessary because
the rational product $B$ below vanishes
at all of them.  Choose five generic
degree-one factors and a degree-five
smoothing section on the smooth
generic ambient; by Bertini and
degree-zero weak Lefschetz they
satisfy all the geometric PG18
conditions.  Spread those data and
choose a rational $q$ unramified
in $E$ outside their finite bad
set and the reference bad set.

Put $B=\prod_{p\in S_0}p$,
including the prime $2$, and
impose $F_1\equiv P+qQ\pmod{q^2}$
and the chosen smooth-section
congruences at the other selected
prime ideals.  Their solution
lattice meets the fixed-reference
Lefschetz open.  After that
pencil is fixed, spread it and
choose integers $a\equiv1\pmod{q^2}$,
$a\equiv1\pmod B$, and $a\equiv0$
at the remaining rational bad
primes.  The forms $q^2BF_0+aF_1$
give the exact product equation
$P+qQ_a$ with $\bar Q_a=\bar Q$
at $q$, smooth reductions at
the other exceptional places,
and good or ordinary quadratic
reductions elsewhere.  The
nonzero special Cartier equations
give flat integral projective
models, exactly as checked in
PA18; the connectedness argument
above handles their special
one-node fibers.

At $q$, apply the uniform surface
curve gate and PG18, then RI18
for every finite extension.
At the other places apply the
frozen good-reduction theorem
or ON17, both with all their
new-cycle extension conclusions.
All finite places and complementary
degrees now satisfy PH16; GA17
supplies the admissible lifts on
this full height domain.
The polynomial-on-each-line
argument of CG18 proves density
in precisely the supplied linear
system.  Finally five degree-one
factors have $\deg_Y(h^5)=10\deg(A^2)$
top points; the extreme weight
calculation proves
\eqref{r18sc:global-monodromy}.
\end{proof}

\paragraph{Boundaries and literature.}
The new ample-cut propagation uses geometric
weak Lefschetz and Poincare duality in their
explicit degree-two ranges; it proves the
needed Gysin injectivity before using it.
The curve property on projective bundles
over surfaces and its smooth-modification
stability were already frozen in CD17.
PG18 supplies numerical local exactness;
CL17/DQ17/ER17/VT17 retain their precise
ordinary-actual reconstruction scope.
RI18, ON17, PH16 and GA17 are used only
where their stated extension and global
hypotheses have been checked.

A local good surface alone does not prove
the global theorem's hypothesis: its
surface must have the supplied smooth
projective model at every finite place.
The separate SG18 example constructs
such an input and checks its literature
dependencies.  The present section
neither proves the general finite-field
curve comparison nor the general G
conjecture or arbitrary-variety full
unconditional heights.
